% ID: SEL-D0022-OB-E04
% ESTADO: cotejada. Segunda pasada visual completa y borrador compilado revisados.
% Fuente: Selectividad/matematicas_aplicadas/reserva_a_EXAMEN 5 -ANDALUCÍA.pdf
% SHA-256: 668689642ca997cf817c9f9f6608d2835efe898e3162b050d7b22b51f225d363
% Página: 2; opción B; ejercicio original 4.

\textbf{(2.5 puntos)} En un artículo de internet se afirma que el número medio de mensajes de WhatsApp que mandan los jóvenes al día no es inferior a 40.

Para contrastar dicha información se elige una muestra aleatoria de 100 jóvenes y se observa que envían una media de 38 mensajes al día. Se sabe que el número de mensajes enviados diariamente sigue una distribución Normal de desviación típica 2. Con un nivel de significación del 5\% plantee un contraste, $(H_0:\mu\geq 40)$, determine la región de rechazo y concluya si \textquestiondown{}se puede aceptar la afirmación del artículo de internet?
